\documentclass[12pt]{article}
\usepackage[utf8]{inputenc}
\usepackage{amsmath,amssymb,amsthm}
\usepackage{algorithm}
\usepackage{algpseudocode}
\usepackage{booktabs}
\usepackage{geometry}
\geometry{margin=1in}

\newtheorem{theorem}{Theorem}
\newtheorem{lemma}[theorem]{Lemma}
\newtheorem{proposition}[theorem]{Proposition}

\title{Problem 114: Counting Block Combinations I}
\author{Project Euler Solution}
\date{}

\begin{document}
\maketitle

\section{Problem Statement}

A row of length $n$ is filled with red blocks (minimum length~3) and black unit squares, such that any two red blocks are separated by at least one black square. How many valid fillings exist for $n = 50$?

\section{Mathematical Development}

\begin{theorem}[Tiling recurrence]
Let $f(n)$ denote the number of valid fillings. Then
\[
f(n) = f(n-1) + \sum_{L=3}^{n} f(n - L - 1)
\]
with $f(-1) = 1$ and $f(0) = 1$.
\end{theorem}

\begin{proof}
The leftmost cell is either black (contributing $f(n-1)$) or the start of a red block of length $L \ge 3$, followed by a mandatory black separator (unless $L = n$). The remaining subproblem has length $n - L - 1$, with $f(-1) = 1$ handling the boundary case $L = n$. The cases are mutually exclusive and exhaustive.
\end{proof}

\begin{lemma}[Prefix-sum optimization]
Define $P(k) = \sum_{j=-1}^{k} f(j)$. Then for $n \ge 3$:
\[
f(n) = f(n-1) + P(n-4).
\]
\end{lemma}

\begin{proof}
Substituting $j = n - L - 1$ transforms $\sum_{L=3}^{n} f(n-L-1)$ into $\sum_{j=-1}^{n-4} f(j) = P(n-4)$. The prefix sum satisfies $P(k) = P(k-1) + f(k)$, so each step costs $O(1)$.
\end{proof}

\begin{theorem}[Verification]
$f(7) = 17$, confirming the problem statement.
\end{theorem}

\begin{proof}
Direct computation:
\begin{center}
\begin{tabular}{rll}
\toprule
$n$ & $f(n)$ & $P(n)$ \\
\midrule
$-1$ & 1 & 1 \\
0 & 1 & 2 \\
1 & 1 & 3 \\
2 & 1 & 4 \\
3 & 2 & 6 \\
4 & 4 & 10 \\
5 & 7 & 17 \\
6 & 11 & 28 \\
7 & 17 & 45 \\
\bottomrule
\end{tabular}
\end{center}
\end{proof}

\section{Algorithm}

\begin{algorithm}[H]
\caption{Compute $f(n)$ via prefix sums}
\begin{algorithmic}[1]
\State $f[-1] \gets 1$;\; $f[0] \gets 1$;\; $P[-1] \gets 1$;\; $P[0] \gets 2$
\For{$i = 1$ to $n$}
    \State $s \gets P[i-4]$ if $i \ge 3$, else $0$
    \State $f[i] \gets f[i-1] + s$
    \State $P[i] \gets P[i-1] + f[i]$
\EndFor
\State \Return $f[n]$
\end{algorithmic}
\end{algorithm}

\section{Complexity Analysis}

\begin{itemize}
\item \textbf{Time:} $O(n)$. Each iteration performs $O(1)$ arithmetic.
\item \textbf{Space:} $O(n)$ for the arrays $f$ and $P$.
\end{itemize}

\section{Answer}

\[
\boxed{16475640049}
\]

\end{document}
